% GENERATED FILE. Edit canonical lessons, then run npm run generate:books.
% canonical-source-hash: fnv1a64:bc260da990e5755b
% canonical-lessons: JA-C263-kyuukyuusha, JA-C263-yobu, JA-C263-sakebu, JA-C263-awateru, JA-C263-sugu

\chapter{Calling for Help}
\label{ch:ja-calling-for-help}
\hlchaptermodality{263}

\begin{chapteropening}
\textbf{By the end of this chapter:} \emph{I can say who called for help and what they did.}

Complete the last lesson of chapter 263: I can say who called for help and what they did.
\end{chapteropening}

\section[\texorpdfstring{\ja{きゅうきゅうしゃ} (kyūkyūsha)}{きゅうきゅうしゃ (kyūkyūsha)}]{\ja{きゅうきゅうしゃ} (kyūkyūsha) — an ambulance}

\label{lesson:JA-C263-kyuukyuusha}



\begin{quote}
Before the new one: say the Japanese for suddenly, then the Japanese for damage.
\end{quote}

\subsection*{You'll want to know: \ja{きゅうきゅうしゃ}}
\textbf{\ja{きゅうきゅうしゃ}} — \emph{kyūkyūsha} — \textquotedblleft{}an ambulance\textquotedblright{}.

\begin{quote}
\textbf{\ja{じこが} \ja{あって、きゅうきゅうしゃが} \ja{きました。}} — \emph{jiko ga atte, kyūkyūsha ga kimashita} — \textquotedblleft{}There was an accident, and an ambulance came.\textquotedblright{}
\end{quote}

\begin{cousinweb}
\textbf{\ja{きゅうきゅうしゃ}} (\emph{kyūkyūsha}) is Sino-Japanese (kango): three Chinese characters read with their on'yomi, readings borrowed from Chinese. They are \emph{kyū}, \textquotedblleft{}rescue\textquotedblright{} (Middle Chinese \emph{kjuwH}), \emph{kyū}, \textquotedblleft{}emergency\textquotedblright{} (\emph{kip}) and \emph{sha}, \textquotedblleft{}car\textquotedblright{} (\emph{tsyhae}) (KANJIDIC2; Middle Chinese in Baxter's transcription).
\end{cousinweb}

\begin{grammarlens}[title={What you've built}]
One more word for telling what happened.
\end{grammarlens}

\subsection*{Your turn}
\begin{itemize}
\raggedright
  \item \emph{Say it:} \emph{kyūkyūsha}
  \item \emph{Say it:} \emph{kyūkyūsha}, once more
  \item \emph{Say it:} \emph{jiko ga atte, kyūkyūsha ga kimashita} — There was an accident, and an ambulance came.
  \item \emph{Recall:} say the Japanese for suddenly, then the Japanese for damage, then say \emph{kyūkyūsha} again
\end{itemize}

\begin{tcolorbox}[breakable,colback=teal!4,colframe=teal!35!black,title={Before you move on}]
What does \ja{きゅうきゅうしゃ} mean? (\textquotedblleft{}An ambulance\textquotedblright{}.) Say it once more.
\end{tcolorbox}
\section[\texorpdfstring{\ja{よぶ} (yobu)}{よぶ (yobu)}]{\ja{よぶ} (yobu) — to call; to send for}

\label{lesson:JA-C263-yobu}



\begin{quote}
Before the new one: say the Japanese for damage, then the Japanese for an ambulance.
\end{quote}

\subsection*{You'll want to know: \ja{よぶ}}
\textbf{\ja{よぶ}} — \emph{yobu} — \textquotedblleft{}to call; to send for\textquotedblright{}.

\begin{quote}
\textbf{\ja{ちちが} \ja{きゅうきゅうしゃを} \ja{よびました。}} — \emph{chichi ga kyūkyūsha o yobimashita} — \textquotedblleft{}My father phoned for an ambulance.\textquotedblright{}
\end{quote}

Polite \textbf{\ja{よびます}}, \emph{yobimasu}; a godan verb. Past (polite): \textbf{\ja{ははが} \ja{いしゃを} \ja{よびました。}}, \emph{haha ga isha o yobimashita}.

\begin{cousinweb}
\textbf{\ja{よぶ}} (\emph{yobu}) is a native Japanese word (wago), inherited rather than borrowed from Chinese; its kanji was attached later as a kun'yomi (JMdict).
\end{cousinweb}

\begin{grammarlens}[title={What you've built}]
One more word for telling what happened.
\end{grammarlens}

\subsection*{Your turn}
\begin{itemize}
\raggedright
  \item \emph{Say it:} \emph{yobu}
  \item \emph{Say it:} \emph{yobu}, once more
  \item \emph{Say it:} \emph{chichi ga kyūkyūsha o yobimashita} — My father phoned for an ambulance.
  \item \emph{Recall:} say the Japanese for damage, then the Japanese for an ambulance, then say \emph{yobu} again
\end{itemize}

\begin{tcolorbox}[breakable,colback=teal!4,colframe=teal!35!black,title={Before you move on}]
What does \ja{よぶ} mean? (\textquotedblleft{}To call; to send for\textquotedblright{}.) Say it once more.
\end{tcolorbox}
\section[\texorpdfstring{\ja{さけぶ} (sakebu)}{さけぶ (sakebu)}]{\ja{さけぶ} (sakebu) — to shout; to cry out}

\label{lesson:JA-C263-sakebu}



\begin{quote}
Before the new one: say the Japanese for an ambulance, then the Japanese for to call.
\end{quote}

\subsection*{You'll want to know: \ja{さけぶ}}
\textbf{\ja{さけぶ}} — \emph{sakebu} — \textquotedblleft{}to shout; to cry out\textquotedblright{}.

\begin{quote}
\textbf{\ja{こどもが} \ja{かわで} \ja{さけびました。}} — \emph{kodomo ga kawa de sakebimashita} — \textquotedblleft{}A child cried out by the river.\textquotedblright{}
\end{quote}

Polite \textbf{\ja{さけびます}}, \emph{sakebimasu}; a godan verb. Past (polite): \textbf{\ja{おとこが} \ja{さけびました。}}, \emph{otoko ga sakebimashita}.

\begin{cousinweb}
\textbf{\ja{さけぶ}} (\emph{sakebu}) is a native Japanese word (wago), inherited rather than borrowed from Chinese; its kanji was attached later as a kun'yomi (JMdict).
\end{cousinweb}

\begin{grammarlens}[title={What you've built}]
One more word for telling what happened.
\end{grammarlens}

\subsection*{Your turn}
\begin{itemize}
\raggedright
  \item \emph{Say it:} \emph{sakebu}
  \item \emph{Say it:} \emph{sakebu}, once more
  \item \emph{Say it:} \emph{kodomo ga kawa de sakebimashita} — A child cried out by the river.
  \item \emph{Recall:} say the Japanese for an ambulance, then the Japanese for to call, then say \emph{sakebu} again
\end{itemize}

\begin{tcolorbox}[breakable,colback=teal!4,colframe=teal!35!black,title={Before you move on}]
What does \ja{さけぶ} mean? (\textquotedblleft{}To shout; to cry out\textquotedblright{}.) Say it once more.
\end{tcolorbox}
\section[\texorpdfstring{\ja{あわてる} (awateru)}{あわてる (awateru)}]{\ja{あわてる} (awateru) — to panic; to be flustered}

\label{lesson:JA-C263-awateru}



\begin{quote}
Before the new one: say the Japanese for to call, then the Japanese for to shout.
\end{quote}

\subsection*{You'll want to know: \ja{あわてる}}
\textbf{\ja{あわてる}} — \emph{awateru} — \textquotedblleft{}to panic; to be flustered\textquotedblright{}.

\begin{quote}
\textbf{\ja{あわてたら、あぶないでしょう。}} — \emph{awatetara, abunai deshō} — \textquotedblleft{}If you panicked, it would be dangerous.\textquotedblright{}
\end{quote}

Polite \textbf{\ja{あわてます}}, \emph{awatemasu}; an ichidan verb. Past (polite): \textbf{\ja{ちちが} \ja{あわてました。}}, \emph{chichi ga awatemashita}.

\begin{cousinweb}
\textbf{\ja{あわてる}} (\emph{awateru}) is a native Japanese word (wago), inherited rather than borrowed from Chinese; its kanji was attached later as a kun'yomi (JMdict).
\end{cousinweb}

\begin{grammarlens}[title={What you've built}]
One more word for telling what happened.
\end{grammarlens}

\subsection*{Your turn}
\begin{itemize}
\raggedright
  \item \emph{Say it:} \emph{awateru}
  \item \emph{Say it:} \emph{awateru}, once more
  \item \emph{Say it:} \emph{awatetara, abunai deshō} — If you panicked, it would be dangerous.
  \item \emph{Recall:} say the Japanese for to call, then the Japanese for to shout, then say \emph{awateru} again
\end{itemize}

\begin{tcolorbox}[breakable,colback=teal!4,colframe=teal!35!black,title={Before you move on}]
What does \ja{あわてる} mean? (\textquotedblleft{}To panic; to be flustered\textquotedblright{}.) Say it once more.
\end{tcolorbox}
\section[\texorpdfstring{\ja{すぐ} (sugu)}{すぐ (sugu)}]{\ja{すぐ} (sugu) — at once; right away}

\label{lesson:JA-C263-sugu}



\begin{quote}
Before the new one: say the Japanese for to shout, then the Japanese for to panic.
\end{quote}

\subsection*{You'll want to know: \ja{すぐ}}
\textbf{\ja{すぐ}} — \emph{sugu} — \textquotedblleft{}at once; right away\textquotedblright{}.

\begin{quote}
\textbf{\ja{けいさつが} \ja{すぐ} \ja{きました。}} — \emph{keisatsu ga sugu kimashita} — \textquotedblleft{}The police came straight away.\textquotedblright{}
\end{quote}

\begin{cousinweb}
\textbf{\ja{すぐ}} (\emph{sugu}) is a native Japanese word (wago), inherited rather than borrowed from Chinese; its kanji was attached later as a kun'yomi (JMdict).
\end{cousinweb}

\begin{grammarlens}[title={What you've built}]
One more word for telling what happened.
\end{grammarlens}

\subsection*{Your turn}
\begin{itemize}
\raggedright
  \item \emph{Say it:} \emph{sugu}
  \item \emph{Say it:} \emph{sugu}, once more
  \item \emph{Say it:} \emph{keisatsu ga sugu kimashita} — The police came straight away.
  \item \emph{Recall:} say the Japanese for to shout, then the Japanese for to panic, then say \emph{sugu} again
\end{itemize}

\begin{tcolorbox}[breakable,colback=teal!4,colframe=teal!35!black,title={Before you move on}]
What does \ja{すぐ} mean? (\textquotedblleft{}At once; right away\textquotedblright{}.) Say it once more.
\end{tcolorbox}
